\begin{helper}
Fix a history cell \(H\), a target event \(T\), and the terminal coordinate
set \(U\).  Assume \(0\le p\le 1/2\) and \(u_B\le u_R\).  Suppose the red
fixed-count event is
covered by a finite family of cells indexed by:
\[
B\in\mathcal B,\qquad J\in\mathcal J,\qquad
\beta\in\mathcal T,\qquad \tau\in\mathcal P .
\]
Here \(\mathcal B\) is a finite set of blue support labels, each a
\(u_B\)-set of blue terminal coordinates in \(U\), and
\(\mathcal J\) is a finite abstract set of red labels.  The index
\(\mathcal T\) is a finite branch/trace label type.  For each
\(\beta\in\mathcal T\) and \(J\in\mathcal J\), there is a decoded red support
\[
R_{\beta,J}\subseteq U.
\]
Whenever the cell indexed by \((B,J,\beta,\tau)\) contains a sample, this
decoded support has cardinality \(u_R\) and consists of red terminal
coordinates.  Assume
\[
|\mathcal J|\le \binom{N_R}{u_R},\qquad
|\mathcal B|\le \binom{N_B}{u_B}.
\]
A transcript \(\tau\) consists of three finite terminal-coordinate sets
\[
P_\tau,\qquad A_\tau,\qquad Z_\tau .
\]
The cell \(G_{B,J,\beta,\tau}\) fixes every coordinate of \(P_\tau\) present,
fixes every coordinate of \(A_\tau\) absent, has
\[
B\cup R_{\beta,J}\subseteq P_\tau,\qquad
Z_\tau\subseteq A_\tau,\qquad
P_\tau\cap A_\tau=\emptyset,
\]
and satisfies
\[
|Z_\tau|\ge \max(0,a-u_R-u_B).
\]
Assume each cell has nonnegative weight \(w_{B,J,\beta,\tau}\), bounded by
the terminal cylinder factor
\[
w_{B,J,\beta,\tau}
\le p^{|P_\tau|}(1-p)^{|A_\tau|},
\]
and that the raw cell mass satisfies
\[
\Pr(T\cap H\cap G_{B,J,\beta,\tau})
\le
\Pr(H)\,w_{B,J,\beta,\tau}.
\]
Finally assume the finite branch/transcript overcount estimate: for every
fixed pair of support labels \(B,J\),
\[
\sum_{\beta,\tau} w_{B,J,\beta,\tau}
\le
p^{u_R}p^{u_B}(1-p)^{\max(0,a-u_R-u_B)}.
\]
Then
\[
\Pr(T\mid H)\le
\binom{N_R}{u_R}p^{u_R}
\binom{N_B}{u_B}p^{u_B}
(1-p)^{\max(0,a-u_R-u_B)} .
\]
\end{helper}
\begin{proof}
Let
\[
L=\max(0,a-u_R-u_B).
\]
The hypotheses give a genuine finite cover of \(T\cap H\): every sample in
\(T\cap H\) belongs to at least one cell \(G_{B,J,\beta,\tau}\), with
\(B\in\mathcal B\), \(J\in\mathcal J\), \(\beta\in\mathcal T\), and
\(\tau\) one of the prefix transcripts.  Since sample weights are
nonnegative, the raw probability of the covered event is bounded by the sum
of the raw probabilities of the cells.  This is exactly the finite
subadditivity mechanism formalized by
\(\noderef{FixedSetHistoryCellBranchTranscriptSummation}\), applied after
combining the two support labels into the support index and using the
branch/trace label and prefix transcript as the branch/transcript indices.
Thus
\[
\Pr(T\cap H)\le
\sum_{B\in\mathcal B}\sum_{J\in\mathcal J}
\sum_{\beta,\tau}
\Pr(T\cap H\cap G_{B,J,\beta,\tau}).
\]

For each cell, the transcript data \(P_\tau,A_\tau,Z_\tau\) record the
terminal values fixed by the branch/trace label and by the prefix scan, not
only the mandatory support and selected-absence coordinates.  The inclusion
\(B\cup R_{\beta,J}\subseteq P_\tau\), the inclusion \(Z_\tau\subseteq A_\tau\), and
the disjointness \(P_\tau\cap A_\tau=\emptyset\) make this a single terminal
cylinder.  Applying \(\noderef{FixedSetHistoryCellFixedCylinderCharge}\) with
the terminal product law for the history cell gives the displayed raw cell
bound
\[
\Pr(T\cap H\cap G_{B,J,\beta,\tau})
\le
\Pr(H)\,w_{B,J,\beta,\tau},
\]
where \(w_{B,J,\beta,\tau}\le p^{|P_\tau|}(1-p)^{|A_\tau|}\).  The complete
cell-geometry condition ensures that every value fixed by the branch label
or transcript is charged through the same present/absent terminal cylinder,
so repeated appearances of the same coordinate do not produce repeated
Bernoulli factors.

Fix \(B\) and \(J\).  The overcount estimate in the statement is the
paper-specific residual summation step: when all branch/trace labels and all
prefix transcripts compatible with these support labels are summed, every
auxiliary terminal value fixed only for the trace or transcript is paired
with the complementary value over the same finite residual coordinate set, so
its total contribution is \(p+(1-p)=1\).  The only factors that remain
uniformly are the decoded mandatory red support \(R_{\beta,J}\), the
mandatory blue support \(B\), and the selected absent coordinates
\(Z_\tau\).  Since \(|R_{\beta,J}|=u_R\) in every nonempty cell,
\(|B|=u_B\), and \(|Z_\tau|\ge L\), and since \(0\le p\le1/2\) gives
\(0\le1-p\le1\), the sum of cell weights for this fixed support pair is at
most
\[
p^{u_R}p^{u_B}(1-p)^L .
\]
This is not a count of traces or transcripts; it is the finite
branch/transcript weight sum after all of their terminal values have been
included in \(P_\tau\) or \(A_\tau\).

Substituting the cell estimates and then the fixed-support overcount estimate
gives
\[
\Pr(T\cap H)\le
\Pr(H)
|\mathcal J|\,|\mathcal B|\,
p^{u_R}p^{u_B}(1-p)^L .
\]
The support-label cardinality assumptions give
\[
|\mathcal J|\le\binom{N_R}{u_R},
\qquad
|\mathcal B|\le\binom{N_B}{u_B}.
\]
Therefore
\[
\Pr(T\cap H)\le
\Pr(H)\binom{N_R}{u_R}p^{u_R}
\binom{N_B}{u_B}p^{u_B}(1-p)^L .
\]
If \(\Pr(H)=0\), the definition of
\(\noderef{TwoBiteConditionalProbability}\) makes \(\Pr(T\mid H)=0\), and
the displayed upper bound is nonnegative because all cell weights are
nonnegative and the support sums are nonnegative.  If \(\Pr(H)>0\), divide
the raw inequality by \(\Pr(H)\).  This yields the claimed conditional
probability bound.
\end{proof}
